% =============== Limitations & Future Work ====================================
% Goal: 0.7-1 page · 4-6 limitation paragraphs
% =============================================================================

We highlight \textbf{seven} known limitations of nautilus-compass v0.7.1
and the empirical evaluation in this paper. The first is the most
consequential and should be read carefully.

\paragraph{No demonstrated behavior steering.}
\label{sec:limitations-behavior}
The empirical contribution of this paper is \emph{drift detection
accuracy} (does our system correctly classify a prompt). The
production claim users actually care about is \emph{behavior
steering} (does Claude, when given the drift alert in its context,
actually avoid the failure mode). These are separate empirical
questions. We have designed an A/B evaluation framework
(\texttt{tests/eval\_behavior\_ab.py}) that pairs each deviation
prompt with two responses (with and without nautilus-compass injection)
judged on four behavioral axes by an independent LLM. Implementation
of the full A/B run requires Anthropic API access for Claude calls
and is left to future work. Until those numbers are reported,
readers should not assume our reported AUC implies improved
production safety. This is the most important caveat in the paper.

\paragraph{Train-on-test concern in the in-set ablation.}
The four-step ablation
(Table~\ref{tab:drift-evolution}) reaches AUC 0.9232 on a
100-prompt set partially constructed by the authors. Step~4
explicitly incorporates patterns from misclassified test prompts
into the anchor set, which by classical interpretation contaminates
the test signal. We mitigate via the held-out test set in
Section~\ref{sec:eval-holdout}, but the in-set 0.9232 should be
read as ``illustrating which design decisions matter'' rather than
as a generalization estimate.

\paragraph{Synthetic drift test set.} Our 100-prompt drift evaluation
set is hand-authored to cover canonical aligned and deviation
patterns. AUC 0.92 reflects performance on this distribution; real
production prompt distributions may differ. We have not yet conducted
a longitudinal study on real user traces, in part because the user
base of nautilus-compass is currently the authors and a small private
deployment.

\paragraph{Single-session-user retrieval ceiling.} On LongMemEval-S
single-session-user questions, even with bge-reranker-v2-m3 we
achieve MRR 0.522. The remaining gap is not closable by any
embedding-based reranker: the questions ask for a specific factual
claim (``What degree did I graduate with?'') buried in a 50-session
chatty haystack. Closing this gap requires LLM-based reasoning over
candidates, which is out of scope for our local-only deployment but
trivial to add when API access is acceptable.

\paragraph{Subset 12 is not the full benchmark.} Our LongMemEval-S
results are on a balanced subset of 12 questions (2 per type). Per-type
variance within the full 500 may yield different aggregate metrics.
Running the full 500 takes approximately 75 minutes on m3 + CPU; we
plan to publish the full-set numbers in the GitHub repository before
journal-track submission.

\paragraph{Cross-domain anchor profiles need expert review.} The
legal, medical, and finance starter anchor sets we ship are written
by software engineers reasoning about each domain's typical patterns,
not by domain experts. Production deployment in any of these
sensitive areas should involve review by qualified domain
practitioners before relying on drift detection signals.

\paragraph{Daemon caches a single anchor profile.} The current
implementation cold-loads one anchor profile at daemon startup;
switching profile requires daemon restart. Multi-profile in-memory
caching with per-request profile selection is planned for v0.8.

\paragraph{Windows m3 cold-load failure modes.} On Windows with
Python 3.14, the Hugging Face Hub client (\texttt{httpx} backend)
intermittently fails to download large models like bge-m3 (2.27~GB),
showing zero-byte progress. We work around this via the ModelScope
mirror but consider it a deployment-environment fragility rather
than a fundamental fix. The recommended production setup is WSL2 or
a Linux/macOS host.

\paragraph{Future directions.} Beyond the items above, we are
particularly interested in: (a) cross-anchor similarity analysis to
detect anchor redundancy and over-coverage, (b) integration with
white-box methods when activation access is available, (c)
automated anchor generation from user history (current implementation
is a regex-pattern v0.1; LLM-based v0.2 is on the roadmap), and
(d) cross-language transfer of anchor sets via multilingual
embedders.
